\chapter{Dinero y medidas}\label{ch:g2:measure}

Reglas, balanzas de cocina, jarras, monedas, relojes y
calendarios: este capítulo trata de los números de la vida diaria ---
cuánto mide, cuánto pesa, cuánto cabe, cuánto cuesta y qué
hora es.

\section{Longitudes, masas y capacidades}

\begin{definition}[Unidades de cada día]\label{def:g2:measure:units}
\begin{itemize}
  \item \emph{Longitud}: el centímetro (cm) para las cosas pequeñas y el metro (m)
        para las grandes, con $1$ m $= 100$ cm;
  \item \emph{masa}\index{masa}: el gramo (g) para lo ligero y el kilogramo (kg) para
        lo pesado, con $1$ kg $= 1000$ g;
  \item \emph{capacidad}\index{capacidad} (lo que cabe en un recipiente): el litro (L).
\end{itemize}
Las unidades más pequeñas y sus reglas exactas llegan en el \cref{ch:g3:measure}.
\end{definition}

\begin{example}[¿Qué unidad?]\label{ex:g2:measure:choose}
Una cera: unos $9$ cm. Una puerta: unos $2$ m. Una manzana: unos $150$
g. Un saco grande de harina: $1$ kg. Una botella de leche: $1$ L. Adivina la unidad
antes de medir --- así se evitan respuestas absurdas como «la puerta mide $2$
cm».
\end{example}

\section{El dinero}

\begin{method}[Formar una cantidad]\label{met:g2:measure:money}
Con monedas y billetes de $1$, $2$, $5$, $10$ y $20$:
\begin{enumerate}
  \item empieza por el mayor que quepa y después completa el resto;
  \item para formar $27$: un $20$, un $5$ y un $2$;
  \item a menudo hay varias maneras --- $27$ también es $10 + 10 + 5 +
        2$.
\end{enumerate}
\end{method}

\begin{omfigure}
\begin{tikzpicture}[scale=0.9]
  \draw[thick, fill=omProp!20] (0,0) rectangle (1.3,0.7);
  \node[font=\small] at (0.65,0.35) {$20$};
  \node[font=\large] at (1.9,0.35) {$+$};
  \draw[thick, fill=omMeth!20] (2.85,0.35) circle (0.35);
  \node[font=\small] at (2.85,0.35) {$5$};
  \node[font=\large] at (3.8,0.35) {$+$};
  \draw[thick, fill=omDef!20] (4.7,0.35) circle (0.3);
  \node[font=\small] at (4.7,0.35) {$2$};
  \node[font=\large] at (5.6,0.35) {$=$};
  \node[font=\normalsize] at (6.5,0.35) {$27$};
\end{tikzpicture}

{\small Formar $27$ con un billete y dos monedas --- y hay otras
maneras, como $10 + 10 + 5 + 2$.}
\end{omfigure}

\begin{example}[El cambio]\label{ex:g2:measure:change}
Un libro cuesta $14$ y pagas con un $20$. El cambio es $20 - 14 =
6$, hallado contando hacia arriba: de $14 \to 20$ hay $6$
(\cref{met:g2:subtraction:countup}).
\end{example}

\section{El tiempo}

\begin{method}[Medias y cuartos de hora]\label{met:g2:measure:clock}
La aguja larga dice los minutos pasados de la hora:
\begin{enumerate}
  \item recta hacia arriba (el $12$): en punto --- $3{:}00$;
  \item un cuarto de vuelta a la derecha (el $3$): y cuarto --- $3{:}15$;
  \item recta hacia abajo (el $6$): y media --- $3{:}30$;
  \item tres cuartos de vuelta (el $9$): menos cuarto para la hora siguiente ---
        $3{:}45$, «las cuatro menos cuarto».
\end{enumerate}
\end{method}

\begin{omfigure}
\begin{tikzpicture}[scale=0.8]
  \foreach \x/\hh/\mm/\lbl in {0/3/3/{3:15}, 4/3/6/{3:30},
                                8/3/9/{3:45}} {
    \begin{scope}[xshift=\x cm]
    \draw[thick] (0,0) circle (1.2);
    \foreach \h in {1,...,12} {
      \node[font=\tiny] at ({90-\h*30}:0.98) {\h};
    }
    \draw[very thick, omDef] (0,0) -- ({90-(\hh+\mm/12)*30}:0.55);
    \draw[thick, omThm] (0,0) -- ({90-\mm*30}:0.95);
    \node[below, font=\small] at (0,-1.45) {\lbl};
    \end{scope}
  }
\end{tikzpicture}

{\small Y cuarto, y media, menos cuarto: la aguja larga roja en
el $3$, en el $6$ y en el $9$.}
\end{omfigure}

\begin{example}[El calendario]\label{ex:g2:measure:calendar}
Un año tiene $12$ meses, en este orden: enero, febrero, marzo, abril,
mayo, junio, julio, agosto, septiembre, octubre, noviembre y diciembre.
Casi todos los meses tienen $30$ o $31$ días; febrero tiene $28$ (o $29$).
Una fecha nombra el día y el mes: «el $5$ de marzo».
\end{example}

\section{Ejercicios}

\begin{exercise}[$\star$]\label{exo:g2:measure:1}
Elige la unidad (cm, m, g, kg o L): un lápiz; el largo de la
clase; una pluma; una mochila; una botella de agua.
\end{exercise}

\begin{exercise}[$\star$]\label{exo:g2:measure:2}
Convierte: $2$ m en cm; $300$ cm en m; $1$ kg en g.
\end{exercise}

\begin{exercise}[$\star$]\label{exo:g2:measure:3}
Forma estas cantidades con monedas y billetes: $16$; \quad $23$; \quad
$35$. Da una manera para cada una.
\end{exercise}

\begin{exercise}[$\star$]\label{exo:g2:measure:4}
Calcula el cambio: pagas $10$ por un juguete de $6$; pagas $20$ por un libro
de $17$; pagas $50$ por un juego de $32$.
\end{exercise}

\begin{exercise}[$\star$]\label{exo:g2:measure:5}
Lee la hora: aguja larga en el $6$ y corta entre el $4$ y el $5$;
aguja larga en el $3$ y corta justo pasado el $8$.
\end{exercise}

\begin{exercise}[$\star$]\label{exo:g2:measure:6}
Dibuja un reloj que marque las dos y media. ¿Adónde apunta cada aguja?
\end{exercise}

\begin{exercise}[$\star$]\label{exo:g2:measure:7}
Nombra el mes: después de mayo; antes de enero; el mes de tu
cumpleaños.
\end{exercise}

\begin{exercise}[$\star$]\label{exo:g2:measure:8}
Lea compra una bebida por $3$ y un pastel por $4$. Paga con un billete de
$10$. ¿Cuánto le devuelven?
\end{exercise}

\begin{exercise}[$\star$]\label{exo:g2:measure:9}
¿Qué pesa más: una bolsa de azúcar de $1$ kg o $800$ g de manzanas? ¿Cuántos
gramos más?
\end{exercise}

\begin{exercise}[$\star\star$]\label{exo:g2:measure:10}
Forma $40$ con exactamente cuatro monedas o billetes (elegidos entre $1$, $2$,
$5$, $10$ y $20$). Encuentra dos maneras distintas.
\end{exercise}
